\chapter{Kidney Transplant}

\section*{What Is This Surgery?}
Kidney transplantation is a definitive surgical procedure involving the implantation of a healthy, functional kidney, sourced from either a deceased or a living donor, into a recipient afflicted with end-stage kidney disease (ESKD). This complex operation is designed to restore critical renal functions, including the filtration of metabolic waste products from the blood, maintenance of fluid and electrolyte homeostasis, and endocrine functions such as erythropoietin production. The donor kidney is typically placed heterotopically within the recipient's iliac fossa, where it is surgically connected to the major iliac blood vessels and the urinary bladder. Successful transplantation offers a profound improvement in the recipient's quality of life, significantly extends survival, and liberates them from the demanding and often debilitating regimen of chronic dialysis.

\section*{Alternative Names}
\begin{itemize}
  \item Renal Transplant
  \item Renal Transplantation
  \item Kidney Graft Surgery
  \item Allogeneic Kidney Transplant
  \item Kidney Replacement Surgery
  \item Cadaveric Kidney Transplant (for deceased donor)
\end{itemize}

\section*{Relevant Surgical Anatomy}
The primary anatomical site for kidney transplantation is the retroperitoneal space of the right or left iliac fossa. This location is strategically chosen due to its proximity to large, easily accessible blood vessels and the urinary bladder, minimizing the length of vascular and ureteral anastomoses. Key anatomical structures involved include the external iliac artery and vein, which serve as the primary recipient vessels for the donor renal artery and vein, respectively. These vessels lie anterior to the psoas major muscle and are readily identifiable. The internal iliac artery and vein are also in close proximity, and occasionally the internal iliac artery may be used for arterial anastomosis, particularly in pediatric recipients or when the external iliac artery is unsuitable.

The urinary bladder is the recipient organ for the donor ureter, typically accessed through its anterior wall. The lymphatic vessels in the iliac fossa, particularly those surrounding the iliac vessels, require careful ligation during dissection to prevent postoperative lymphocele formation. The native kidneys are generally left in situ unless specific indications for nephrectomy exist (e.g., uncontrolled hypertension, recurrent infections, polycystic kidney disease causing mass effect). The peritoneum is reflected medially to expose the retroperitoneal space, ensuring the transplanted kidney lies extraperitoneally, which reduces the risk of intra-abdominal complications.

\section{Surgical landmarks and considerations:}
\begin{itemize}
  \item \textbf{External Iliac Artery and Vein:} These are the primary targets for vascular anastomoses, providing robust flow.
  \item \textbf{Urinary Bladder:} The dome of the bladder is used for ureteroneocystostomy.
  \item \textbf{Retroperitoneal Space:} The ideal location for graft placement, minimizing adhesions and infection risk.
  \item \textbf{Lymphatics:} Careful ligation is crucial to prevent lymphocele.
  \item \textbf{Psoas Muscle:} Provides a soft bed for the transplanted kidney.
\end{itemize}

\section*{Surgical Principle \& Rationale}
The fundamental surgical principle of kidney transplantation is the heterotopic placement of a functional renal allograft to replace the excretory and endocrine functions of the failed native kidneys. The rationale is to restore physiological homeostasis by providing effective glomerular filtration, tubular reabsorption, and endocrine activity, thereby reversing the uremic state. The donor kidney is typically positioned in the iliac fossa, a site that offers several advantages: it is extraperitoneal, minimizing the risk of bowel injury and peritonitis; it provides direct access to large, high-flow recipient vessels (external iliac artery and vein); and it allows for a relatively short, tension-free ureteral anastomosis to the bladder.

The technical goals of the surgery are paramount: first, to establish robust, patent, and tension-free vascular anastomoses between the donor renal artery and vein and the recipient's iliac vessels, ensuring immediate and adequate blood flow to the graft. Second, to create a watertight, non-obstructive ureteroneocystostomy between the donor ureter and the recipient's bladder, facilitating unimpeded urine drainage and preventing urine leakage or stricture. Physiologically, a successful transplant aims for rapid reperfusion of the graft, leading to immediate or early diuresis and a progressive decline in serum creatinine, signifying the restoration of renal function. Lifelong immunosuppression is a critical adjunct to the surgical principle, preventing the recipient's immune system from recognizing and rejecting the foreign donor organ.

\section*{History \& Surgical Evolution}
The history of kidney transplantation is a testament to perseverance and scientific innovation. Early attempts at organ transplantation in the early 20th century, notably by Alexis Carrel, laid foundational surgical techniques but were hampered by the then-ununderstood immunological barriers. The true breakthrough occurred on December 23, 1954, when Dr. Joseph Murray and his team at Peter Bent Brigham Hospital in Boston performed the first successful human kidney transplant between identical twin brothers, Richard and Ronald Herrick. This landmark procedure bypassed the immune rejection problem due to their genetic identity, earning Dr. Murray the Nobel Prize in Physiology or Medicine in 1990.

The subsequent decades witnessed a rapid evolution driven by the development of effective immunosuppressive agents. The 1960s saw the introduction of azathioprine and corticosteroids, which offered the first pharmacological means to suppress the immune response in non-identical donor-recipient pairs, albeit with significant side effects and limited success rates. The landscape of transplantation was revolutionized in the early 1980s with the advent of cyclosporine, a calcineurin inhibitor. This drug dramatically improved graft survival rates by specifically targeting T-cell activation, transforming kidney transplantation from an experimental procedure into the preferred treatment for ESKD. Further advancements in immunosuppression, including tacrolimus, mycophenolate mofetil, and sirolimus, alongside refined surgical techniques, improved organ preservation, and enhanced patient management, have continued to expand eligibility criteria and optimize long-term outcomes, making kidney transplantation a routine and highly successful surgical intervention.

\section{Clinical Indications}
Kidney transplantation is indicated for individuals diagnosed with end-stage kidney disease (ESKD) who meet specific medical, surgical, and psychosocial criteria. The primary goal is to restore renal function and improve quality of life.

\begin{itemize}
  \item To treat end-stage kidney disease (ESKD) resulting from various etiologies, including diabetes mellitus, hypertension, glomerulonephritis, polycystic kidney disease, and other chronic nephropathies.
  \item To improve patient survival and quality of life for individuals currently undergoing chronic dialysis (hemodialysis or peritoneal dialysis), as transplantation offers superior outcomes compared to long-term dialysis.
  \item To provide pre-emptive treatment for progressive chronic kidney disease (CKD) that is rapidly approaching ESKD, allowing transplantation before the initiation of dialysis, which is associated with better graft and patient outcomes.
  \item To address kidney failure in pediatric patients, where transplantation is crucial for promoting normal growth, development, and long-term psychosocial well-being, offering superior outcomes to dialysis.
  \item To resolve severe uremic symptoms that are refractory to dialysis, such as intractable pruritus, peripheral neuropathy, uremic pericarditis, or severe malnutrition.
  \item To manage specific metabolic disorders (e.g., primary hyperoxaluria) that primarily affect the kidneys and can be definitively treated by replacing the diseased organ.
  \item To offer a more cost-effective long-term treatment option compared to lifelong dialysis, considering the overall healthcare burden and improved productivity of transplant recipients.
  \item To treat patients with certain forms of atypical hemolytic uremic syndrome (aHUS) or C3 glomerulopathy where kidney replacement is necessary.
\end{itemize}

\section{Clinical Positioning}
Kidney transplantation is widely recognized as the primary and most definitive treatment modality for the majority of patients with end-stage kidney disease (ESKD). Its clinical positioning as a first-line therapy is supported by overwhelming evidence demonstrating superior patient survival, graft survival, and overall quality of life compared to chronic dialysis. For eligible patients, particularly those without significant comorbidities, a pre-emptive transplant, performed before the initiation of dialysis, represents the ideal scenario. This approach avoids the complications, lifestyle burdens, and cardiovascular risks associated with dialysis, leading to better long-term outcomes.

However, in certain clinical contexts, kidney transplantation may function as a secondary or salvage procedure. This applies to patients who have experienced a failed previous kidney transplant and require re-transplantation. Similarly, individuals who have been on dialysis for an extended period due to factors such as delayed listing on the transplant waitlist, donor availability issues, or temporary medical contraindications, will receive a transplant as a secondary intervention after their initial life-sustaining management with dialysis. While dialysis serves as a crucial bridge therapy, transplantation offers the potential for a return to a more normal physiological state and a significantly improved prognosis, making it the ultimate therapeutic goal for most ESKD patients.

\section{Surgical Technique \& Procedure}
Kidney transplantation is a meticulously performed surgical procedure requiring general anesthesia, typically lasting between 3 to 5 hours. The patient is positioned supine on the operating table, and the surgical field is prepped and draped in a sterile fashion.

\begin{itemize}
  \item \textbf{Anesthesia and Incision:} General anesthesia is induced, and the patient is positioned. A curvilinear incision, commonly known as a "hockey stick" or Gibson incision, is made in the right or left lower quadrant of the abdomen. This incision extends from the pubic symphysis superiorly and laterally towards the anterior superior iliac spine, providing optimal exposure to the retroperitoneal space, iliac vessels, and bladder.
  \item \textbf{Vascular Dissection:} The abdominal wall muscles are carefully divided, and the peritoneum is reflected medially to expose the retroperitoneal space. The external iliac artery and vein are meticulously identified, dissected free from surrounding tissues, and prepared for anastomosis. Lymphatic vessels in the area are carefully ligated or clipped to minimize the risk of postoperative lymphocele formation, a common complication.
  \item \textbf{Kidney Placement and Vascular Anastomoses:} The donor kidney, which has been preserved in a cold storage solution (e.g., University of Wisconsin solution) to minimize ischemic injury, is brought into the sterile field. It is typically placed in the iliac fossa. The donor renal vein is anastomosed end-to-side to the recipient's external iliac vein using fine vascular sutures (e.g., 6-0 or 7-0 Prolene). Subsequently, the donor renal artery (or arteries, if multiple) is anastomosed end-to-side to the recipient's external iliac artery. Meticulous surgical technique is paramount to ensure patent, tension-free anastomoses, crucial for immediate graft function.
  \item \textbf{Reperfusion and Assessment:} Once the vascular anastomoses are completed, vascular clamps are carefully removed, allowing blood flow to be restored to the transplanted kidney. The kidney is observed for immediate reperfusion, indicated by a rapid change in color from pale to pink, and often, immediate urine production from the donor ureter. Any bleeding from the anastomotic sites is meticulously controlled.
  \item \textbf{Ureteral Anastomosis (Ureteroneocystostomy):} A small incision is made in the dome of the recipient's bladder. The donor ureter is then brought through a tunnel created in the bladder wall (to prevent reflux) and anastomosed to the bladder mucosa, typically using a Lich-Gregoir or Politano-Leadbetter technique. A ureteral stent (e.g., a double J stent) is often placed temporarily to ensure patency, prevent stricture, and minimize the risk of urine leakage during the initial healing phase. The stent is usually removed endoscopically several weeks post-transplant.
  \item \textbf{Closure:} After confirming meticulous hemostasis, satisfactory graft perfusion, and urine output, a surgical drain (e.g., a Jackson-Pratt drain) may be placed in the retroperitoneal space to monitor for fluid collections. The abdominal wall layers are then closed meticulously with absorbable sutures, and the skin is closed with staples or non-absorbable sutures.
\end{itemize}

\section*{Preoperative Workup \& Preparation}
Extensive and multidisciplinary preoperative workup and preparation are critical for optimizing recipient selection, minimizing risks, and ensuring the long-term success of kidney transplantation. This comprehensive evaluation aims to identify and address any medical conditions that could compromise the transplant or the patient's survival.

\begin{itemize}
  \item \textbf{Comprehensive Medical Evaluation:} A thorough assessment of cardiac function (including ECG, echocardiogram, and often stress testing or coronary angiography), pulmonary function (PFTs), gastrointestinal health, and dental clearance is performed to rule out occult infections or significant comorbidities.
  \item \textbf{Infectious Disease Screening:} Extensive screening for active and latent infections (e.g., HIV, Hepatitis B and C, CMV, EBV, tuberculosis, syphilis, toxoplasmosis) is mandatory. Any active infections must be treated and resolved prior to transplantation to prevent severe post-transplant complications under immunosuppression.
  \item \textbf{Immunological Workup:} ABO blood group compatibility and Human Leukocyte Antigen (HLA) tissue typing are performed for both donor and recipient to assess immunological match. A crossmatch test, which detects pre-formed antibodies in the recipient against donor antigens, is critical to prevent hyperacute rejection.
  \item \textbf{Psychosocial Evaluation:} A comprehensive psychosocial assessment ensures the recipient has adequate social support systems, possesses a clear understanding of the lifelong commitment to immunosuppression, and demonstrates the capacity for strict adherence to complex medical regimens.
  \item \textbf{Cancer Screening:} Age-appropriate cancer screening (e.g., colonoscopy, mammography, prostate-specific antigen, dermatological evaluation) is completed to ensure no active malignancy exists, as immunosuppression can accelerate cancer growth or recurrence.
  \item \textbf{Medication Adjustments:} Patients receive detailed instructions on which medications to hold (e.g., anticoagulants, certain antihypertensives) and which to continue prior to surgery. Dialysis may be performed shortly before surgery to optimize fluid and electrolyte balance and ensure the patient is euvolemic.
  \item \textbf{NPO Status:} The patient must be nil per os (NPO) for a specified period, typically 6-8 hours, before surgery to minimize the risk of aspiration during general anesthesia.
  \item \textbf{Pre-operative Antibiotics:} Prophylactic broad-spectrum antibiotics are administered intravenously shortly before the surgical incision to reduce the risk of surgical site infection and other opportunistic infections.
  \item \textbf{Informed Consent:} A detailed discussion of the surgical procedure, its potential benefits, risks, alternatives, and the lifelong implications of immunosuppression is conducted, and comprehensive informed consent is obtained from the patient.
\end{itemize}

\section*{Contraindications \& Precautions}
Careful consideration of contraindications and precautions is essential for patient safety and graft longevity in kidney transplantation.

\textbf{Absolute Contraindications:}
\begin{itemize}
  \item Active, uncontrolled systemic infection or sepsis, which would be exacerbated by immunosuppression.
  \item Active malignancy with metastatic disease or a high risk of recurrence within 2-5 years, as immunosuppression can accelerate cancer progression.
  \item Severe, irreversible cardiac disease (e.g., severe cardiomyopathy, intractable angina, severe valvular disease) that precludes safe surgery or long-term survival.
  \item Active substance abuse (alcohol or illicit drugs) or uncorrected severe psychiatric illness leading to documented non-adherence to medical therapy.
  \item Untreated or uncontrolled severe peripheral vascular disease or cerebrovascular disease that poses an unacceptable surgical risk.
  \item Documented and persistent non-adherence to medical therapy or follow-up, indicating a high risk of graft failure due to non-compliance with immunosuppression.
  \item Irreversible brain damage or other conditions incompatible with a reasonable quality of life or functional recovery post-transplant.
\end{itemize}

\textbf{Relative Contraindications \& Precautions:}
\begin{itemize}
  \item Advanced age (typically >70-75 years), which may increase surgical and immunosuppression-related risks, requiring careful patient selection and individualized risk assessment.
  \item Morbid obesity (BMI >35-40 kg/m\textasciicircum{}2), increasing surgical complexity, wound complications, delayed graft function, and potentially impacting long-term outcomes.
  \item Significant vascular calcification of the iliac vessels, which can complicate vascular anastomoses and increase the risk of thrombosis or hemorrhage.
  \item Certain autoimmune diseases (e.g., lupus nephritis, focal segmental glomerulosclerosis) with a high risk of recurrence in the transplanted graft, requiring careful pre-transplant management and specific post-transplant protocols.
  \item Significant pelvic anatomical abnormalities or prior extensive pelvic surgery, which may make surgical access and anastomoses technically challenging.
  \item Active peptic ulcer disease or diverticulitis, which may be exacerbated by corticosteroid use and immunosuppression, requiring pre-transplant treatment.
  \item HIV infection, which is now considered a relative contraindication with effective antiretroviral therapy and careful monitoring, but still requires specialized management.
  \item Chronic liver disease or other significant organ dysfunction that may limit overall survival or complicate immunosuppression management.
\end{itemize}

\section*{Complications \& Risks}
Kidney transplantation, while life-saving, is a major surgical procedure carrying a spectrum of potential risks, both immediate and long-term, related to the surgery itself and the lifelong requirement for immunosuppression.

\textbf{General Surgical Complications:}
\begin{itemize}
  \item \textbf{Bleeding:} Intraoperative or postoperative hemorrhage, potentially requiring blood transfusions or re-operation.
  \item \textbf{Infection:} Surgical site infection (SSI), wound dehiscence, or deep space infection (e.g., abscess).
  \item \textbf{Anesthesia Complications:} Adverse reactions to anesthetic agents, respiratory or cardiac complications (e.g., myocardial infarction, arrhythmia), stroke, or, rarely, death.
  \item \textbf{Thromboembolic Events:} Deep vein thrombosis (DVT) and pulmonary embolism (PE), despite prophylactic measures.
  \item \textbf{Nerve Injury:} Damage to peripheral nerves (e.g., lateral femoral cutaneous nerve) during incision or dissection, leading to numbness or pain.
\end{itemize}

\textbf{Procedure-Specific Risks:}
\begin{itemize}
  \item \textbf{Graft Thrombosis:} Acute clotting of the renal artery or vein, leading to immediate graft failure, often necessitating emergency re-operation or graft removal.
  \item \textbf{Delayed Graft Function (DGF):} The transplanted kidney does not function immediately, requiring temporary dialysis. This is often due to acute tubular necrosis (ATN) from ischemia-reperfusion injury.
  \item \textbf{Ureteral Complications:} Urine leak from the ureteral anastomosis, ureteral stricture, or obstruction, potentially requiring stent placement, percutaneous nephrostomy, or surgical revision.
  \item \textbf{Lymphocele:} Collection of lymphatic fluid around the graft, which may cause compression of the ureter or vessels, requiring percutaneous drainage or surgical marsupialization.
  \item \textbf{Renal Artery Stenosis:} Narrowing of the renal artery anastomosis, leading to new-onset or worsening hypertension and graft dysfunction, often managed with angioplasty and stenting.
  \item \textbf{Graft Rejection:} The recipient's immune system attacks the donor kidney. This can range from hyperacute (immediate, rare with modern crossmatching) to acute (days to months post-transplant) or chronic (gradual decline over years).
  \item \textbf{Graft Rupture:} A rare but catastrophic complication, often associated with acute rejection, leading to severe hemorrhage.
  \item \textbf{Recurrence of Native Disease:} The original kidney disease (e.g., FSGS, IgA nephropathy) can recur in the transplanted kidney.
\end{itemize}

\textbf{Immunosuppression-Related Risks (Long-Term):}
\begin{itemize}
  \item \textbf{Increased Susceptibility to Infections:} Opportunistic infections (e.g., Cytomegalovirus (CMV), Epstein-Barr Virus (EBV), Pneumocystis jirovecii pneumonia (PJP), fungal infections, urinary tract infections) due to a weakened immune system.
  \item \textbf{Cardiovascular Disease:} Accelerated atherosclerosis, hypertension, dyslipidemia, and new-onset diabetes after transplant (NODAT) are common side effects of immunosuppressants, increasing the risk of heart attack and stroke.
  \item \textbf{Malignancy:} Increased risk of certain cancers, including skin cancers (squamous cell carcinoma, melanoma), post-transplant lymphoproliferative disorder (PTLD), Kaposi's sarcoma, and cervical cancer.
  \item \textbf{Bone Disease:} Osteoporosis, osteopenia, and avascular necrosis, particularly associated with long-term corticosteroid use.
  \item \textbf{Nephrotoxicity:} Some immunosuppressants (e.g., calcineurin inhibitors like tacrolimus and cyclosporine) can be toxic to the transplanted kidney or native kidneys, contributing to chronic graft dysfunction.
  \item \textbf{Other Side Effects:} Tremor, gastrointestinal upset, hirsutism, gingival hyperplasia, mood changes, and electrolyte disturbances.
\end{itemize}

\section*{Postoperative Recovery Timeline}
The postoperative recovery timeline for kidney transplant recipients is a carefully managed process, beginning immediately after surgery and extending over months to years.

Immediately after kidney transplantation, the recipient is transferred to the Post-Anesthesia Care Unit (PACU) for intensive monitoring of vital signs, urine output, and pain control. Once stable, they are typically moved to a specialized transplant ward or an intensive care unit (ICU) for the first 24-48 hours. During this critical period, continuous, often hourly, monitoring of urine output is paramount, alongside frequent assessment of serum creatinine, blood urea nitrogen (BUN), and electrolyte levels to gauge initial graft function. Intravenous fluids are meticulously managed to maintain adequate hydration and support the new kidney. Immunosuppressive medications, typically a multi-drug regimen including a calcineurin inhibitor (e.g., tacrolimus), an antiproliferative agent (e.g., mycophenolate mofetil), and corticosteroids, are initiated immediately to prevent acute rejection. Early ambulation is encouraged to prevent thromboembolic complications.

Over the first 1-2 weeks of hospitalization, the focus shifts to optimizing immunosuppression drug levels, managing pain effectively, preventing infections with prophylactic antibiotics and antivirals, and promoting progressive ambulation and self-care. Patients are educated extensively on their complex medication regimen, potential side effects, and warning signs of complications. Dietary restrictions are gradually liberalized, and physical therapy assists with mobility. The ureteral stent, if placed, is usually removed endoscopically several weeks post-transplant. Long-term recovery involves lifelong immunosuppression and regular monitoring. Most patients can return to light activities within 4-6 weeks and resume normal activities, including work, within 2-3 months, though heavy lifting is restricted for several months to allow abdominal wall healing. Full recovery and adaptation to the new lifestyle often take 6-12 months.

\section*{Surgical Findings \& Pathology}
During kidney transplantation, the surgeon meticulously documents several key intraoperative findings. These include the condition of the recipient's iliac vessels (e.g., presence of atherosclerosis, calcification), the ease of dissection in the retroperitoneal space, the appearance of the donor kidney upon reperfusion (e.g., immediate pinking, turgor, urine production), the quality of the vascular anastomoses (e.g., absence of leaks, good thrill), and the integrity of the ureteroneocystostomy. Any anatomical variations or unexpected findings are also noted. If native nephrectomy is performed concurrently, the gross appearance of the diseased native kidneys (e.g., polycystic changes, severe hydronephrosis, tumor) is described.

Postoperatively, the pathology report on any resected tissues, such as native kidneys, provides a definitive diagnosis of the underlying end-stage kidney disease. More commonly, if graft dysfunction occurs, a percutaneous kidney biopsy of the transplanted kidney is performed. The pathology report from this biopsy is crucial for diagnosing the cause of dysfunction, differentiating between various types of rejection (e.g., T-cell mediated, antibody-mediated), acute tubular necrosis (ATN), drug toxicity, or recurrence of the original disease. It details specific histological findings, such as inflammatory infiltrates, vascular changes, glomerular injury, or fibrosis, guiding subsequent therapeutic interventions and informing the long-term prognosis of the graft.

\section{Expected Postoperative Course}
A normal and successful postoperative course following kidney transplantation is characterized by a progressive return to health and independence. Patients typically experience a significant reduction in uremic symptoms, such as fatigue, nausea, and pruritus, often within days to weeks of the surgery. Pain at the incision site is managed with analgesics and gradually subsides over the first few weeks. The most important indicator of success is the rapid and sustained improvement in kidney function, reflected by a progressive decrease in serum creatinine and blood urea nitrogen (BUN) levels to a normal or near-normal range.

Patients should experience robust urine output, without signs of obstruction or leakage. As kidney function normalizes, fluid and electrolyte imbalances resolve, and the need for dialysis ceases. Energy levels improve, appetite returns, and patients can gradually resume normal daily activities. The surgical wound heals without complication, and any temporary ureteral stent is removed. A successful course means the transplanted kidney is effectively sustaining life, allowing the recipient to return to a more active and fulfilling life, free from the constraints of dialysis, while adhering to their lifelong immunosuppression regimen.

\section{Postoperative Care \& Follow-up}
Lifelong postoperative care and rigorous follow-up are absolutely essential after kidney transplantation to monitor graft function, manage immunosuppression, and detect complications early.

\begin{itemize}
  \item \textbf{Frequent Clinic Visits:} Initially, weekly or bi-weekly visits for the first few months, gradually decreasing to monthly, then quarterly, and eventually annually, depending on the patient's stability and graft function.
  \item \textbf{Laboratory Tests:} Regular blood tests including serum creatinine, BUN, electrolytes, complete blood count (CBC), liver function tests (LFTs), and therapeutic drug monitoring of immunosuppressant levels (e.g., tacrolimus, cyclosporine) to ensure optimal ranges and monitor for side effects.
  \item \textbf{Urine Tests:} Urinalysis and urine protein-to-creatinine ratio to screen for proteinuria, hematuria, or signs of infection.
  \item \textbf{Imaging Studies:} Periodic Doppler ultrasound of the renal graft to assess blood flow, rule out hydronephrosis, and monitor for fluid collections (e.g., lymphocele).
  \item \textbf{Kidney Biopsy:} Performed if there is a rise in creatinine or other signs of graft dysfunction, to diagnose rejection or other pathologies.
  \item \textbf{Blood Pressure Monitoring:} Close monitoring and aggressive management of blood pressure, as hypertension is common post-transplant and can impact graft survival.
  \item \textbf{Diabetes Screening:} Regular screening for new-onset diabetes after transplant (NODAT), a common side effect of immunosuppression.
  \item \textbf{Infection Surveillance:} Monitoring for opportunistic infections, with appropriate vaccinations and prophylactic medications (e.g., for CMV, PJP).
  \item \textbf{Cancer Screening:} Enhanced and regular screening for skin cancers, post-transplant lymphoproliferative disorder (PTLD), and other malignancies due to increased risk with immunosuppression.
  \item \textbf{Bone Health Monitoring:} Assessment for osteoporosis and management with calcium, vitamin D, and bisphosphonates if needed.
  \item \textbf{Medication Adherence Counseling:} Ongoing education and support to ensure strict, lifelong adherence to the complex immunosuppression regimen, which is critical for graft survival.
  \item \textbf{Suture/Staple Removal:} Typically occurs 10-14 days post-surgery.
  \item \textbf{Rehabilitation/Physical Therapy:} Encouraged for gradual return to activity and strength.
\end{itemize}
Patients are instructed to immediately contact their transplant team for any fever, decreased urine output, new or worsening swelling, or new abdominal pain.

\section*{Epidemiology}
Kidney transplantation stands as the most frequently performed solid organ transplant globally, a reflection of the high and increasing prevalence of end-stage kidney disease (ESKD). In the United States alone, over 25,000 kidney transplants are performed annually, yet more than 100,000 individuals remain on the waiting list, highlighting a significant disparity between demand and donor organ availability. The global incidence of ESKD is steadily rising, primarily driven by the escalating rates of chronic diseases such as diabetes mellitus and hypertension, which collectively account for the majority of new ESKD cases. While ESKD can affect all age groups, its highest incidence is typically observed in individuals over 60 years of age. There is a slight male predominance in the prevalence of ESKD and, consequently, among kidney transplant recipients. The mortality rate for patients on the transplant waiting list is substantial, underscoring the critical and life-saving nature of transplantation. Post-transplant, patient and graft survival rates have shown continuous improvement over the decades, though long-term morbidity remains a concern due to the lifelong need for immunosuppression and its associated complications, including infection, cardiovascular disease, and malignancy.

\section*{Outcomes \& Prognosis}
The outcomes of kidney transplantation have dramatically improved over the past decades, establishing it as the most successful form of organ replacement therapy. For deceased donor transplants, 1-year graft survival rates typically range from 90-95\%, with 5-year graft survival rates around 75-85\%. Living donor transplants generally yield even superior results, boasting 1-year graft survival rates exceeding 95\% and 5-year rates often above 85-90\%. Patient survival rates are also significantly higher for transplant recipients compared to patients remaining on dialysis, with a 5-year survival advantage of 10-20\% or more, depending on recipient characteristics. Functional outcomes are excellent, including a return to normal daily activities, improved energy levels, resolution of uremic symptoms, and a substantial enhancement in overall quality of life.

Recurrence of the original kidney disease in the transplanted graft is a concern for certain conditions like focal segmental glomerulosclerosis (FSGS), IgA nephropathy, or atypical hemolytic uremic syndrome, but this risk is managed with specific pre- and post-transplant protocols. Prognostic factors significantly influencing long-term outcomes include the type of donor (living vs. deceased), the degree of HLA matching, cold ischemia time of the donor organ, recipient age and comorbidities, and critically, the patient's adherence to the complex lifelong immunosuppressive therapy. Landmark clinical guidelines, such as those from the Kidney Disease: Improving Global Outcomes (KDIGO) organization, consistently emphasize the superior long-term survival and quality of life benefits of kidney transplantation over dialysis for eligible patients.

\section{References}
\begin{enumerate}
  \item MedlinePlus. Kidney transplant. U.S. National Library of Medicine; 2024. Available from: \url{https://medlineplus.gov/kidneytransplant.html}
  \item Hricik DE, et al. Kidney Transplantation. In: Brenner and Rector's The Kidney. 11th ed. Elsevier; 2020.
  \item National Kidney Foundation. Kidney Transplant. Available from: \url{https://www.kidney.org/atoz/content/transplantinfo}
  \item Kidney Disease: Improving Global Outcomes (KDIGO) Transplant Work Group. KDIGO Clinical Practice Guideline for the Care of Kidney Transplant Recipients. Kidney International Supplements. 2009;1(2):S1-S103.
\end{enumerate}

\clearpage